Der fallende Ball erreicht die halbe Höhe zur Zeit $t$ mit
\begin{align}
 \frac{h}{2} &= \frac{1}{2}g\,t^2 \quad\Rightarrow\quad t = \sqrt{\frac{h}{g}}
\end{align}
Der geworfene Ball muss zur selben Zeit auf derselben Höhe sein:
\begin{align}
 \frac{h}{2} &= v_0\,t - \frac{1}{2}g\,t^2
\end{align}
Addiert man die beiden Gleichungen, fallen die Terme mit $g$ heraus: $h = v_0\,t$. Also
\begin{align}
 v_0 &= \frac{h}{t} = \result{\sqrt{g\,h}}
\end{align}
Im Treffpunkt ist der geworfene Ball übrigens gerade zur Ruhe gekommen: $v_0 - g\,t = \sqrt{gh}-\sqrt{gh} = 0$.
